\section{The periodic vacuum and its physical pairing}\label{c6:periodic-polynomial}

Polynomial exchange equations, degree bounds and special-parameter fusion have useful precedents in loop models. Di Francesco and Zinn-Justin~\cite{DiFrancescoZinnJustin2005} use a multiparameter polynomial construction in the dense $O(1)$ model. Garbali and Nienhuis~\cite{GarbaliNienhuis2017Ground,GarbaliNienhuis2017Sum} develop polynomial $q$KZ vectors for the dilute $O(1)$ model with open boundaries and derive a determinant normalization from a first fusion recurrence. Here the periodic $n=0$ vectors, their degrees and their physical pairing are proved with the normalization required for the two marked observables.


Here is a transfer calculation we will use on the honeycomb cylinder. Keep the parameters $L,b,S,\kappa,s$ and the local operators $R(u),Q,W,\Omega$ as in \eqref{eq:legacy-P-2}, \eqref{eq:legacy-P-3} (\(s=e^{2iL}\), \(\kappa=(2\cos L)^{-1}\)); put \(C=2-\sqrt 2=4\cos^2 b\). Let \(N\ge1\) be an integer and \(m,n\ge0\) integers with \(m+n=N\). Use the infinite-height lozenge-cylinder, period \(N\) in the first lattice coordinate, and choose corners \(A,B\) of tiles on the same horizontal lattice line, \(A\) just before column 1, \(B\) just after column \(m\). Let \(\mathcal Z\) sum systems of finitely many disjoint simple polygons in the cylinder (each with weight \(\kappa^{\rm length}\), length in triangle visits). All must separate \(A,B\) on the cylinder, and each gets an additional factor \(y+y^{-1}\). Include the empty system.

The finite formula uses spectral variables \(v_j\); its physical interpretation will be at their common value zero. There are two groups of sites \(1,\ldots,m\) and \(m+1,\ldots,N\). Put
\[
 \begin{gathered}
 R_s(d)=-1/[\sin(d+L)\sin(d-L)],\quad R_c(d)=\cos^2 d,\\
 D(d)=\sin(d+2L)\sin(d-2L)\sin(d+b)\sin(d-b).
 \end{gathered}
\]
Sum over site labels \(\alpha_j=+1,0,-1\), with sums of labels in the two groups equal to one another, say \(h\). Include a factor \(y^h\prod_j c_{\alpha_j}\) with \(c_\pm=\sqrt C,\ c_0=2\), and a two-body factor for each \(i<j\): write \(d=v_j-v_i\), and take \(P_{\alpha_i\alpha_j}(d)\) given by
\[
 \begin{array}{c|cccc}
 (\alpha,\alpha') &(++),(--)&(00)&(+0),(0-)& (+,-)\\ \hline
 P_{\alpha\alpha'}(d)&1&H(d)&H(d-L)&H(d-L)H(d-2L),
 \end{array}\qquad P_{\alpha\alpha'}(d)=P_{\alpha'\alpha}(-d).
\]
Here \(H=R_s\) for a same-group pair and \(R_c\) across groups. For each same-group pair also include \(D(d)\). This sum of products is \(G\), with removable singularities understood.

We define its denominator for both parities of \(N\). Put
\[
 B(u)=\frac{\sin(u+b)\sin(u-b)}{\sin u},\qquad
 h(u)=\frac{\cos u}{B(u)}.
\]
Let \(\mathsf h_N\) be the skew matrix with upper entries
\((\mathsf h_N)_{jk}=h(v_k-v_j)\), \(j<k\). For odd \(N\), append a last column of ones and a row of minus ones, with bottom-right entry zero. Define
\[
 f_N(v)=a_N\prod_{j<k}B(v_k-v_j)\,\Pf\mathsf h_N,
 \qquad a_{2l}=2^l/\sqrt2,\quad a_{2l+1}=(-2)^l.
\]
The empty product and empty Pfaffian are 1, so \(f_0=1/\sqrt2\) and \(f_1=1\). The proof below shows that the apparent singularities are removable and that \(f_N\) is a trigonometric polynomial. For even size this definition gives
\[
 2f_{2m'}^2=4^{m'}\prod_{i<j}B(v_j-v_i)^2
       \det[h(v_j-v_i)]_{i,j=1}^{2m'}.
\]

\begin{proposition}[Homogeneous physical pairing]\label{prop:physical-pairing}
For the cylinder, corners and polygon weights just defined, and for \(y>0\),
\[
 \mathcal Z=\left.\frac{G}{2f_N^2}\right|_{v_1=\cdots=v_N=0}.
\]
The denominator is nonzero there. The equality also holds as a Laurent-polynomial identity in \(y\), with the stated even and odd conventions for \(f_N\).
\end{proposition}

The proof first bounds the denominators of the periodic eigenvectors and then uses their wheel zeros to determine their degrees and fusion rules. These bounds reduce the pairing formula to a polynomial identity fixed by fusions and one balanced residual test; the final orientation calculation identifies its homogeneous value with the physical partition.
\subsection{Periodic eigenlines}\label{c6:subheading-04}

Write \(K(\xi)_{pp}=\xi^p\) for the diagonal spin twist, \(p=+,0,-\). In this calculation transfer is single-row, on site spins, and we use the sector of total spin zero:
\[
 T_\xi(t)=\operatorname{Tr}_a K_a(\xi)\,R_{a1}(t+v_1)\cdots R_{aN}(t+v_N),\qquad \xi=i\ \text{or }-i.
\]
We use the empty site vector (all zeros), denoted by index \(\varnothing\). Site order is as displayed even though operator action is from the right. By trigonometric degree at most \(D\) we mean Laurent exponents between \(-D\) and \(D\) in \(e^{iv_j}\). If the parity under \(v_j\mapsto v_j+\pi\) is \((-1)^D\), multiplication by \(e^{iDv_j}\) gives a polynomial of degree at most \(D\) in \(e^{2iv_j}\); hence \(D+1\) distinct zeros modulo \(\pi\) force it to vanish.

\begin{auxlemma}[Polynomial vacuum]
For \(N\ge1\) and generic rapidities, the family \(T_i(t)\) has a simple common eigenline with eigenvalue 1. Normalize its vector by \(w^i_\varnothing=1\), and put \(\Psi_N=f_Nw^i\), with \(\Psi_0=f_0=1/\sqrt2\). Then \(\Psi_N\) is a trigonometric polynomial of degree at most \(N-1\) in each rapidity. On a component of site spin \(p_j\), its parity in \(v_j\) is \((-1)^{N-1+p_j}\).

For an adjacent pair at the beginning of the indicated order, with remaining list \(U\), its exact fusions are
\[
\begin{aligned}
 \Psi_N(u,u+2L,U)
 &=\prod_{z\in U}\cos(z-u-L)\,
       W_{12}\Psi_{N-1}(u+L,U),\\
 \Psi_N(u,u+b,U)
 &=2\cos b\prod_{z\in U}\cos(z-u-L)\cos(z-u-2L)\,
       \Omega_{12}\Psi_{N-2}(U).
\end{aligned}
\]
At either imaginary infinity of the last rapidity,
\[
 \frac{\Psi_N(v)}{\prod_{j<N}\cos(v_N-v_j)}
 \longrightarrow \sqrt2\,\Psi_{N-1}(v_1,\ldots,v_{N-1})\otimes|0\rangle_N.
\]
\end{auxlemma}

\begin{proof}[Proof of the polynomial vacuum lemma]
We first construct rational normalized eigenvectors for both twists. The denominator bound and wheel calculation will show that \(f_N\) clears the twist-\(i\) vector with exactly the degree in the statement. All equalities in spectral variables are used first at generic values.

The operators at common twist commute by Yang--Baxter and spin conservation (interchange two rows by the invertible auxiliary crossing \(R(t-t')\) under the traces). Also, writing \(T_j=T_\xi(-v_p+jL)\) at any site \(p\),
\begin{equation}
 T_0T_2=T_1,\qquad T_0T_3=I,\qquad T_1T_3=T_2.                 \label{eq:legacy-P-11}
\end{equation}
Indeed in the two traced spaces (second row parameter lower than first by \(2L\)), each site's product preserves the image of \(W\) by \eqref{eq:legacy-P-3}, acting there by the center-parameter row. At site \(p\), the two factors are \(R_{2p}(0)R_{1p}(2L)\) (1 and 2 here label traced spaces), with entire range in that image. The two twists preserve the subspace and restrict there to the same twist. This proves the first identity. The second follows identically with difference \(b\), using \(\Omega\) and \eqref{eq:legacy-P-2}, and the third follows from the first two.

Any finite product at common twist has \(\varnothing,\varnothing\) entry 1. Indeed draw the path expansion as for \(R\) above, with positive arrows right to left and bottom to top, and include the twist for seam crossings (positive crossing for motion towards the left). The tangent-turn phase on each loop in the square drawing is \(\exp(i\,{\rm turn}/4)\), with arc corrections telescoping as before. A contractible simple loop has turn \(\pm2\pi\), an essential simple loop has turn 0 and horizontal winding \(\pm1\), so both cancel between orientations. For instance for the essential turn, start at a highest point and follow many repeated periods of the simple lift, closing above by segments from the highest endpoints; the planar turn bound forces turn per period 0. Simplicity of a noncontractible circle implies primitive winding: its two-way infinite lift divides its containing band into an upper and a lower component by planar separation, so the disjoint component lifts (if several) would be vertically ordered, permuted transitively with order preserved by horizontal period translation. They are finite in number (the absolute winding), hence just one. These arguments do not require positive tile weights. In particular the tuple taking value 1 on every \(T_\xi(t)\) (the vacuum eigenvalue) occurs on total spin zero: the entry gives evaluation at 1 on the commuting operator polynomials, a character. By joint generalized eigenspace decomposition of commuting matrices it must occur, and the corresponding spectral projection \(P^\xi\) has \(P^\xi_{\varnothing,\varnothing}=1\). Indeed, take finitely many ordinary row parameters spanning the family, and decompose into joint generalized eigenspaces (successive ordinary primary decompositions). A generic linear combination separates their tuples from one another and from the proposed character if distinct. Its minimal polynomial must vanish at the character by the entry property, so the character occurs. The projection onto it is a polynomial in this combination taking value 1 there.

This eigenvalue is joint algebraically simple for generic rapidities, and we will need simplicity also at generic points of the pair hyperplanes \(v_j-v_k=rL\bmod\pi,\ r=0,\pm1,\pm2,\pm3\). Here and below joint algebraic simplicity means that a generic finite linear combination at ordinary parameter values has the vacuum-line eigenvalue simple. First, one can increase the number of sites by adding a site with \(v_N\) (or similarly any \(v_j\)) near positive or negative imaginary infinity. For definiteness at positive imaginary infinity, its \(R\) tends to a matrix block triangular with respect to that site's spin \(p\), with diagonal blocks \(K_a(s^{-3p})\), and only a possible off-diagonal transition from \(-1\) to \(+1\); the spin-zero block of that site decouples. This follows from the displayed weights of \(R\) (at negative imaginary infinity all signs reverse). In the decoupled block the smaller operator is unchanged. Neither of the two nonzero-spin blocks can give eigenvalue identically 1: send test \(t\) to \emph{opposite} imaginary infinity in the limiting operator. All the other \(R\)'s are triangular in the auxiliary space in a common direction, and the traced result on the reduced spin sector is scalar \(1+\xi s^{-6p}+\xi^{-1}s^{6p}\ne1\) (reverse exponents in the opposite case). Thus simplicity persists, by taking a linear combination simple at the limit. The normalized projection and vectors also extend there. This gives generic simplicity by induction starting with at most one site.

For the indicated pair hyperplanes start instead with just the two sites. Shifts by \(\pi\) are harmless by \eqref{eq:legacy-P-2}. For positive differences \(r=1,2,3\), sample with a zero argument at the rapidity followed by the other at distance \(rL\) (whichever their slot order), to get on the two sites a copy of \(R(rL)\) multiplied by a one-site twist (slot order or multiplication order immaterial here). On total spin zero this has one eigenvalue 1, trace 1 and trace of the square \(1+2c(rL)^2\), using the entry notation for \(R\) in \eqref{eq:legacy-P-2}; \(|c(rL)|<1\), so 1 is simple. For difference 0 sample at common tile value \(L\). The submatrix deleting the empty site state has diagonals \(v(L)^2=\kappa^4\) (here \(v(\cdot)\) is the straight-tile coefficient) and off-diagonals of absolute value \(\kappa^4\), by direct multiplication (\(d(L)=0\)). Its spectral radius is less than 1. The resolvent entry on the empty state is \(1/(z-1)\) since all power entries are 1; the principal cofactor formula gives simplicity. This proves the assertions by adding ordinary sites as above.

Normalize the resulting right column \(w^\xi=P^\xi|\varnothing\rangle\), left row \(\ell^\xi=\langle\varnothing|P^\xi\), so both empty entries are 1 and \(P^\xi=w^\xi\ell^\xi\). These are rational functions of the variables \(e^{iv_j}\), by the eigen-equations. They are regular at generic points of the just-treated hyperplanes, and at either single-site imaginary infinity they tend to the lower vectors with that site empty. By crossing in \eqref{eq:legacy-P-2},
\[
 \ell^\xi(v)=w^{\xi^{-1}}(-v)^T.
\]
Indeed the site transpose of a row also reverses all rapidities about \(b\) in the tile arguments by crossing, flipping auxiliary spins and telescoping phases \(s^{k-i}\).

Here are exchange and fusion rules for the normalized vectors (permuted orders are written with site labels traveling with rapidities). For consecutive sites \(p,q\), with \(p\) then \(q\), \(u=v_q-v_p\),
\[
 w(\cdots,p,q,\cdots)=R_{pq}(u)w(\cdots,q,p,\cdots),\quad
 w(2,\ldots,N,1)=K_1(\xi)w(1,2,\ldots,N).
\]
Yang--Baxter and moving a slot under the trace give the intertwiners (use spin conservation in moving the twist). To see the scalar normalizations, in any indicated order take the two-nonzero entries for sites \(p',q'\) with \(p'\) to the left of \(q'\), all others empty. They are in the fixed proportion \(s^{-1},s\) for spins \(+-,-+\) respectively. In fact this holds after any finite sequence of common-twist rows on the empty input, hence for the projection: the sole arch at the top, drawn from \(p'\) heading down to \(q'\) heading up (case \(-+\)), has turn \(+\pi\) for a lift exiting to its right, or \(-\pi\) for one exiting to its left, by closing at the top. Its lift endpoint displacement has magnitude less than one period by simplicity and noninterlacing with its translates. Thus phases including seam are \(s\) or \(s^{-1}\xi=\pm s\), respectively, and reversing gives \(s^{-1}\) or \(\pm s^{-1}\) with the same sign. In the reversed-order vector of the exchange formula the two contributions to the empty entry after applying \(R_{pq}\) thus cancel. Cyclic movement plainly preserves the empty entry.

For order \(p,q\), at difference \(+2L\) the right vector is \(W_{pq}\) applied to the lower vector with these sites fused to one at \(v_p+L\); at difference \(b\) it is \(\Omega_{pq}\) times the vector with both deleted. Indeed the exchange formula and regularity give the image restriction, and the row acts there as claimed by \eqref{eq:legacy-P-2}, \eqref{eq:legacy-P-3}. To see the latter slot conventions, swap 1 and 2 and flip all spins in the right fusion formulas (spin flip on a pair interchanges \(R_{ij},R_{ji}\), and the simultaneous transformations preserve the embeddings). The empty entry fixes the scalar. The left row at the same respective differences, multiplied on the right by \(W_{pq}\) or by the insertion of \(\Omega_{pq}\), likewise gives the lower row by intertwining and uniqueness. Its normalization follows from crossing and the two-entry phase rule; the extra terms cancel. All these fusion statements concern generic points of the indicated hyperplanes.

\subsection{Denominator bound}\label{c6:subheading-05}

\paragraph{Possible pole divisors.}
We include details because degree control for these rational vectors is important in interpolating the observable. For each sign of twist let \(F_N^\xi\) clear the minimal common denominator of \(w^\xi\), in centered trigonometric convention. More precisely, pole divisors on the complex torus in \(e^{iv_j}\), with their maximum orders among entries, are unchanged under \(v_j\mapsto v_j+\pi\) at each site (the one-site spin sign conjugates the row). Form the product of the corresponding prime Laurent factors to these maximum orders, normalized by a monomial shift so it is polynomial in \(e^{iv_j}\) with no variable as a factor. It transforms proportionally to itself under the individual sign changes, hence evenly in each variable (otherwise divisible by the variable). Thus this minimal clearing factor is a polynomial in \(z_j=e^{2iv_j}\). The divisors are symmetric by exchange and invertibility, away from \(v_j-v_p=\pm2L,\pm b\bmod\pi\), where there are no divisor poles by the simplicity above. They are invariant under simultaneous translations. Thus this polynomial is symmetric (alternation impossible by regularity at coincidences), has equal degrees \(d_N^\xi\) in the variables, and is homogeneous. Conjugation for real rapidities just flips all spins in \(T\); thus the denominator is invariant as a divisor under \(z_j\mapsto1/\bar z_j\) simultaneously. Its total homogeneous degree is consequently \(N d_N^\xi/2\). Multiply it by \(e^{-i d_N^\xi\sum v_j}\) to get \(F_N^\xi\), translation invariant and of trigonometric degree \(d_N^\xi\) in each variable, of definite parity under shifts by \(\pi\). Also \(F_N^\xi w^\xi\) is trigonometric polynomial of degrees at most \(d_N^\xi\), by regularity of the normalized vector at single-site infinities. All these divisor arguments can equivalently be made by factorization in Laurent polynomial rings.

We claim
\begin{equation}
 d_N^i+d_N^{-i}\le 3(N-1)\qquad(N\ge 2).                           \label{eq:legacy-P-12}
\end{equation}
Use for now \(N\ge3\). On total spin zero the row is \(\pi\)-periodic in \(t\) and tends at both imaginary infinities in \(t\) to identity (triangular auxiliary matrices as above). Thus generically
\[
 T_\xi(t)=I+\sum_k [A_k\cot(t+v_k+2L)+B_k'\cot(t+v_k+3L)],\qquad
 \sum_k(A_k+B_k')=0
\]
with commuting residue operators; primes on \(B_k'\) just distinguish notation. Write their column as \(\mathcal R=(A_1,B'_1,\ldots,A_N,B'_N)^t\). At a fixed site \(p\), let \(c_j\) be the scalar row for which \(X_j=T_\xi(-v_p+jL)-I=c_j\mathcal R\); its two entries at site \(q\) are
\[
 (c_j)_{q,A}=\cot(v_q-v_p+(j+2)L),\qquad
 (c_j)_{q,B}=\cot(v_q-v_p+(j+3)L).
\]
The two fusion equations give commuting-operator rows annihilating \(\mathcal R\):
\[
 \mathcal A_{p,0}=c_0-c_1+c_2+X_0c_2,\qquad
 \mathcal A_{p,1}=c_1-c_2+c_3+X_1c_3.
\]
Indeed these are respectively \(X_0-X_1+X_2+X_0X_2=0\) and \(X_1-X_2+X_3+X_1X_3=0\). On the vacuum line the operator coefficients \(X_j\) vanish. With rows \(\mathcal A_{p,1},\mathcal A_{p,0}\) in site order, the resulting scalar matrix is skew:
\[
 H_{JK}=r(y_K-y_J),\quad y=(v_1,v_1+L,\ldots,v_N,v_N+L),\quad
 r(u)=\cot(u+3L)-\cot(u+4L)+\cot(u+5L).
\]
There is also the sum row \({\bf1}^t\), since \(\sum_k(A_k+B'_k)=0\). For the \(2N\) operator rows, or after replacing any row by this sum row, let \(\mathcal A\) denote the operator matrix and \(H_{\mathcal A}\) its scalar restriction to the vacuum. Then
\[
 \det\mathcal A=(\det H_{\mathcal A})P^\xi.
\]
To see this, apply the adjugate to \(\mathcal A\mathcal R=0\): the determinant kills every residue operator. On each nonvacuum joint primary block some residue has nonzero eigenvalue, since its row eigenvalue is not identically 1. A generic linear combination of the residues is therefore invertible on the sum of these blocks, so \(\det\mathcal A\) vanishes there. On the simple vacuum line its value is \(\det H_{\mathcal A}\). This proves the identity and reduces pole control of \(P^\xi\) to scalar determinants.

Outside the pair hyperplanes \(v_k-v_j=r'L\bmod\pi\) at integer \(r'\), the operator entries here are analytic in rapidities. We analyze the scalar Pfaffian cleared as
\[
 J_N=\operatorname{Pf}(H)\prod_{j<k}p(v_k-v_j),\qquad
 p(u)=\frac{\sin(u+2L)\sin(u-2L)\cos u\, [\sin(u+b)\sin(u-b)]^2}
 {\sin^2u\,\sin(u+L)\sin(u-L)} .
\]
Here the Pfaffian uses upper entries in slot order. Recall that its pairing expansion is the coefficient of the ordered top wedge in \(\omega^N/N!\) for the two-form of the skew matrix. Thus it changes by the determinant under simultaneous row and column transformation and squares to the determinant (reduce a nonsingular skew form to independent paired slots by elimination, then extend by polynomial identity). By this expansion \(\operatorname{Pf}(H) H^{-1}{\bf1}\) likewise has entries given up to sign by Pfaffians with the corresponding node replaced by one with constant coupling 1 (expand along that node, or use Pfaffian row replacement and multiply by \(H\)). Then \(J_N\) is a polynomial of definite trigonometric parity and degree at most \(3(N-1)\) per variable. Indeed coincidence of sites causes two coinciding row pairs, giving a double zero, and difference \(\pm L\) gives a simple zero at least. At difference \(2L\) there is only one singular entry between the two pairs of rows, at \(3L\) at most two poles in a matching, and similarly at negatives. At \(4L\) all entries between those pairs have a residue block proportional to \(\left(\begin{smallmatrix}-1&1\\1&-1\end{smallmatrix}\right)\), so double-pole terms cancel. The bounds at infinity then give the degree.

\paragraph{Simple poles and nondegeneracy.}
Here are useful nondegeneracy details. For two sites of difference \(u\), with \(t_*=r(L)=-1/(1+S)\),
\[
 J_2=p(u)[t_*^2-r(u)^2+r(u+L)r(u-L)]
     =\cos u\, [S t_*-(t_*^2/2)\cos 2u].
\]
Indeed the bracket is regular at \(4L\) (even locally and at most a simple pole), leaving an even polynomial of degree 2 after removing \(\cos u\). Writing \(x=\cos 2u\), \(p(u)/\cos u=x(x+S)^2/[2(1-x)(x-S)]\); at infinity this gives the leading term, and at \(u=2L\) the value follows from the simple pole of \(r(u+L)\). The three roots modulo \(\pi\) are simple. Moreover at each of these roots some component of \(J_2 H^{-1}{\bf1}\) is analytic and nonzero. After removing \(p\), Pfaffian cofactors here have first entries \(r(u)-r(u-L)-t_*\), \(r(u)-r(u+L)+t_*\). At \(4L\) they have nonzero residue. At \(x=2S/t_*=-1-\sqrt2\) they cannot both vanish: \(r(u-L)+r(u+L)-2r(u)\) there is
\[
 \sin(2u)\,[-2/x+6/(x+S)-4/(x+1)]\ne0.
\]

For generic fixed other variables, \(J_N\), as a function of \(v_N\), likewise has simple zeros; at each, outside the special differences, \(J_N H^{-1}{\bf1}\) is nonzero. To verify the generic assertion one can take the other imaginary parts very widely separated. Near each \(v_N=v_j+u\) with \(u\) in a fixed bounded region, the two nodes of every distant site have identical limiting couplings to all other nodes (constants from \(r\) at imaginary infinity). They eliminate as a pair with internal upper entry \(t_*\), leaving the remaining skew matrix unchanged in the limit. This works also for Pfaffian cofactors summed against \({\bf1}\), i.e. Pfaffians with one node omitted (taken in the near cluster) and a constant coupling node appended. Thus, up to common nonzero scalar renormalizations and nonvanishing exponential factors from the \(p\)'s, the Pfaffian expression and the indicated cofactor expressions on near-cluster nodes tend to the two-site expressions. This convergence is locally uniform near their roots including \(4L\), since multiplying by \(\cos u\) removes possible poles there also in the cofactor expressions (same rank-one residue). Small circles around the three two-site roots therefore give three distinct simple roots per near cluster, exhausting the degree bound; the nonvanishing assertion persists. This open-set test suffices for generic pole counting in one variable.

\paragraph{The exceptional difference \(4L\).}
For \(N\ge3\), \(J_N\) at \(v_N-v_j=4L\) is itself generically nonzero. For three sites tending to \(0,4L,e\), consider the pole of the Pfaffian before multiplying by the pair factors. Terms using both matchings between the first two pairs (i.e. contracting their cross block twice) have cancelling pole part by the two-site computation. The remaining residue from a single cross contraction with the rank-one block is up to sign the determinant of couplings from the \emph{sums} of nodes of those two pairs to the last pair. Its rows are
\[
 \begin{pmatrix}
 \cot(e+2L)+\cot(e+5L)&\cot(e+3L)+\cot(e+6L)\\
 \cot(e-2L)+\cot(e+L)&\cot(e-L)+\cot(e+2L)
 \end{pmatrix}.
\]
The determinant is not identically zero (pole at \(e=3L\)). Adding distant sites as above preserves nonvanishing after normalization. At this difference \(4L\), where sampling rows could have singularities, use at each of the two involved sites \(T_0T_3=I\) instead of \(T_1T_3=T_2\), and multiply the retained equation \(X_0-X_1+X_2+X_0X_2=0\) by the vanishing difference \(v_N-v_j-4L\). The replacement operator row is \(c_0+c_3+X_0c_3\); its scalar restriction is the sum of the former two scalar rows. In the retained row use \(X_0\) as the coefficient of \(c_2\). Now the operator rows are analytic there (at these two sites \(X_0,X_3\) are regular and the cotangent rows for \(X_1,X_2\) have at most simple poles). Their scalar determinant is up to sign \((v_N-v_j-4L)^2\det H\), since the replacement tangent row is the sum of the former two. It is nonzero at generic points there. This excludes a pole divisor of \(P^\xi\) at this difference, and at its translates by periodicity.

Now outside the integer-\(L\) differences, the determinant identities show that \(P^\xi\) can have poles only at roots of \(J_N\) and of order at most one along the generic test just used: at a simple Pfaffian zero replace a row by the sum row, choosing a nonzero component of \(J_NH^{-1}{\bf1}\) so the new determinant has just a simple zero. The remaining integer-\(L\) differences besides \(4L\) were treated by simplicity earlier. Thus the common pole denominator of \(P^\xi\) has degree at most \(3(N-1)\) in one \(z\)-variable (each root modulo \(\pi\) counted once). Since its entries are the outer products of entries of \(w^\xi(v), w^{\xi^{-1}}(-v)\), both with empty entry 1, their maximum pole orders are exactly the sums of maximum pole orders of the two vectors. This proves \eqref{eq:legacy-P-12}. For two sites with difference \(u\), use the known phase rule to write the zero-spin column in slot order \(+-,00,-+\) as \((s^{-1}E,1,s E)^T\). The eigen-equation at zero first tile argument (\(R_{12}(u)K_1(\xi)\)) gives
\[
 E=\frac{S\sin u}{Q(u)- (s^2/\xi)\sin(b-u)\sin(2L-u)}.
\]
At \(\xi=i\) its denominator after cancellation is proportional to \(\cos u\); at \(\xi=-i\) it is proportional to \(\cos L+\sin L\cos 2u\) without further cancellation. Thus the degrees are 1 and 2 in this case.

\subsection{Wheel zeros}\label{c6:subheading-06}

Let \(a=1\) for \(\xi=i\), \(a=2\) for \(\xi=-i\). The polynomial column \(\Psi=F_N^\xi w^\xi\), with sites in order \(1,2,3,Y\) (remaining list \(Y\)), vanishes at the wheel \((v_1,v_2,v_3)=(u,u+2L,u+5L)\); for \(\xi=-i\) it vanishes to second order at least. Here the other values may be generic. Indeed perturb \(v_2,v_3\) by \(z,z'\) keeping \(v_1=u\). On \(z=0\) the column belongs to the image of \(W_{12}\). On \(z'=z\) it is a scalar times \(b_0=\Omega_{23} w_{1,Y}\), and on \(z'=0\) it is a scalar times
\[
 c_0=\left(\prod_{j\in Y}^{\rm order}R_{3j}(v_j-v_3)\right) K_3(\xi) J_1^{\rm spin}\Omega_{31}\,w_{2,Y},
 \qquad J^{\rm spin}_{pp}=(-1)^p.
\]
Here the subscripts of \(w\) give the smaller list. Indeed move site 3 through \(Y\), use cyclic covariance and \(v_1-v_3=b-\pi\) for the ordered pair \(3,1\). The scalar multipliers are \(F_N^\xi\), and these smaller columns and intertwiners are regular at the wheel point generically. Modulo the image of \(W_{12}\), \(b_0\) at that point is nonzero, and for twist \(-i\), \(b_0,c_0\) are independent there, generically. One checks by sending the other sites successively to imaginary infinity, decoupling as before. On their empty entries the limiting columns are \(|0\rangle_1\Omega_{23}\) and \(K_3 J_1^{\rm spin}\Omega_{31}|0\rangle_2\). For a combination \(\gamma b_0-\alpha c_0\) of those columns to be in the image, at spin \(p=\pm1\) on site 3 the equal coefficients required at spins \(0,-p\) and \(-p,0\) on sites 1,2 give \(\gamma s^p=\alpha(\xi/s)^p\). At spin 0 the image requirement gives \(\gamma=\alpha\). This proves the assertions. Thus \(\Psi\) vanishes at the intersection, and for twist \(-i\) its derivatives along \(z'=z\) and \(z'=0\) there give \(\gamma b_0,\alpha c_0\) with difference in the image of \(W_{12}\), so both zero. In particular the symmetric scalar \(F_N^\xi=\Psi_\varnothing\) has the stated orders, also at permutations and translates modulo \(\pi\).

Restrict \(F_N^\xi\) to \(v_1=u,\ v_2=u+b\), not an identical zero. It clears the denominator for the remaining list, by vector fusion. Also for each remaining variable \(v_j\) it has wheel factors \([\sin(v_j-u+2L)\sin(v_j-u+3L)]^a\). They are independent of the smaller-list denominator factors. Counting degrees in the third variable for \(N\ge3\) gives \(d_N^\xi\ge d_{N-2}^\xi+2a\). The base degrees are \(d_1^i=d_1^{-i}=0\), \(d_2^i=1\) and \(d_2^{-i}=2\). Induction therefore gives
\[
 d_N^i\ge N-1,\qquad d_N^{-i}\ge2(N-1).
\]
Their sum reaches the upper bound \eqref{eq:legacy-P-12}, so both inequalities are equalities. Temporarily write \(\widetilde f_N=F_N^i\), with scale unspecified. Saturation of degrees in all remaining variables gives the exact scalar fusion, up to a nonzero constant:
\[
 \widetilde f_N(u,u+b,U)\ \propto\ \widetilde f_{N-2}(U)\prod_{z\in U}\cos(z-u-L)\cos(z-u-2L).
\]
Indeed there is no remaining Laurent degree span in those variables after dividing the factors, no monomial shift by the centered degree bounds, and simultaneous translation removes any dependence on \(u\). Similarly at \(v_2=u+2L\) there is the denominator for the list \((u+L,U)\), by injectivity of \(W\), and a wheel factor \(\cos(z-u-L)\) per remaining variable. For \(N\ge4\) the smaller-list denominator does not already contain any such factor (absence of the difference-\(4L\) divisor above). This gives
\[
 \widetilde f_N(u,u+2L,U)\ \propto\ \widetilde f_{N-1}(u+L,U)\prod_{z\in U}\cos(z-u-L).
\]

Set \(\widetilde f_0=1/\sqrt2, \widetilde f_1=1, \widetilde f_2(v_1,v_2)=\sqrt2\cos(v_1-v_2)\). For \(N=3\) symmetry, parity and translation invariance give a constant plus a multiple of \(\sum_{j<k}\cos 2(v_k-v_j)\). Matching the difference-\(b\) fusion fixes their ratio; take
\[
 \widetilde f_3=1/\sqrt2+\tfrac12\sum_{j<k}\cos 2(v_k-v_j).
\]
It gives the stated difference-\(2L\) fusion as well, with coefficient 1 as for \(\widetilde f_2\). Take coefficient 1 in that fusion thereafter by scaling. Then the coefficient in the difference-\(b\) fusion is \(2\cos b\) for all \(N\): this holds for \(2,3\) by the formulas, and for higher \(N\) it is the same as for \(N-1\), by commuting the two fusions on disjoint pairs. Indeed for pairs at \((x,x+b)\) and \((u,u+2L)\) the cross factors in each computation multiply to \(\cos(r-L)\cos(r-2L)\cos(r+L)\cos r,\ r=u-x\).

The Pfaffian expression defining \(f_N\) has these same properties. Coincidence poles from \(B\) are removed by coinciding nodes in the Pfaffian, and poles of \(h\) also cancel. Its degrees, parity and symmetry are as required. At \((u,u+b,U)\) only the paired term survives; the identity
\[
 B(d)B(d-b)=-\cos(d-L)\cos(d-2L)
\]
gives the required \(b\)-fusion with the stated coefficients \(a_N\). Starting with sizes 0 and 1, this fusion and symmetry characterize the polynomial: in one variable the difference has degree at most \(N-1\) with fixed parity and vanishes at \(2(N-1)\) distinct prescribed values modulo \(\pi\). Hence \(f_N=\widetilde f_N\). In particular \(f_N(-v)=f_N(v)\).

At either imaginary infinity in \(v_N\), the same Pfaffian gives
\[
 f_N/\prod_{j<N}\cos(v_N-v_j)\longrightarrow\sqrt2 f_{N-1}.
\]
Indeed \(B/\cos\) and \(h\) in the last-site column tend to \(\pm i\) and \(\mp i\), respectively. For odd size the constant couplings of this last site to the others eliminate against the appended node. Combining this scalar rule with the normalized vector's empty-site limit proves the asserted limit for \(\Psi_N\). The degree bound for \(F_N^iw^i\), the scalar fusions just proved and the normalized vector fusions give all remaining claims of the lemma. Its component parity follows from the one-site spin-sign covariance and the parity \((-1)^{N-1}\) of \(f_N\).
\end{proof}

\subsection{Interpolation of the two-point function}\label{c6:subheading-07}

\begin{proof}[Proof of Proposition~\ref{prop:physical-pairing}]
The polynomial vacuum lemma now supplies the finite-degree space in which to identify the pairing. Crossing gives \(\ell^{-i}(v)=w^i(-v)^T\), so \(f_N^2\) clears its two vector factors. At general rapidities consider
\[
 Z=\ell^{-i}(v)\left(\prod_{j\le m}K_j(\delta)\right) w^i(v),\qquad
 \delta=i y,\qquad U=2 f_N^2 Z.
\]
Thus \(U\) is polynomial of degrees \(\le 2(N-1)\), \(\pi\)-periodic in each variable by spin-sign covariance. It is separately symmetric in the two groups by spin conservation and exchange (the exchanged row of the projection transforms with the inverse crossing, by \eqref{eq:legacy-P-2} and reversed arguments). Within either group, on a pair \((x,x+2L)\) or \((x,x+b)\), \(Z\) reduces to \(Z\) of the smaller list without prefactors (by embedding the right column, passing any spin twists through the embedding, and using left fusion). At imaginary infinity of one variable it also limits to \(Z\) with that variable removed.

We show \(G=U\). First \(G\) too is a separately symmetric, \(\pi\)-periodic trigonometric polynomial. At a same-group difference 0 the possible simple residues cancel under label exchange. At difference \(+L\) the only residues are from \(00\) and \(+,-\), cancelling by opposite residues of \(R_s\) at \(\pm L\), \(C R_s(0)=4\), and the identity
\[
 P_{+a}(u)P_{-a}(u-L)=P_{0a}(u)P_{0a}(u-L)
\]
for either kernel. Other poles are removed by symmetry or by \(D\). For growth of a summand in the \(i\)-th variable, each \(P_{a a'}\) has \(1-aa'\) kernel factors. Counting growth degree \(-2\) per same-group and 2 per opposite-group kernel, with degree 4 per \(D\), gives at most \(2(N-1)-2\alpha_i^2\) by label neutrality. Divided by \(\prod_{j\ne i}\cos^2(v_i-v_j)\), the limit therefore comes from \(\alpha_i=0\); since \(P_{0a}(d)=H(d+a L)\), the leading phases from the shifts cancel by neutrality and this limit is 2 times \(G\) for the remaining list. This is also the rule for \(U\).

Within a group \(G\) has the same fusions as \(2 f_N^2 Z\). Here are details. At the pair \((x,x+2L)\) only \(+0,0-,+,-\) survive and reduce to \(+, -,0\) respectively at the fused site. The common pair multiplier including site coefficients (relative to the fused site) is \(2\lim_{d\to2L} D(d)R_s(d-L)=1\), using \(C R_s(0)=4\). For any surviving pair \(a,a'\) the ratio of external factors to another site, of label \(c'\) at \(x+u\), uses
\[
 P_{a c'}(u)P_{a' c'}(u-2L)/P_{a+a',c'}(u-L)=H(u-L)
\]
directly by the table. Including \(D(u)D(u-2L)/D(u-L)\) for the same-group case this multiplier is \(\cos^2(u-L)\) in both cases. At \((x,x+b)\) only \(+,-\) remains, of net pair multiplier \(C\); the external kernel factors multiply to \(H(u-L)H(u-2L)\). Including \(D(u)D(u-b)\) within-group gives \(\cos^2(u-L)\cos^2(u-2L)\) in both cases. All constants here follow by \(8L=\pi\) in the displayed kernel and sine products.

The initial sizes 0,1 give equality directly. By induction the difference \(G-U\) is divisible by the product of within-group \(D\)'s and has degrees at most \(2(N-1)-2\) by the infinity rule. After division by that product, the degree bounds are
\[
 2(N-1)-2-4(m-1)=2(n-m)\quad\text{on the first group},
\]
with \(2(m-n)\) on the second. One is negative when \(m\ne n\), so the difference vanishes. When \(m=n\), only a constant multiple of the product is possible. Evaluate in this case at
\[
 (x_1,\ldots,x_m,x_m+b,\ldots,x_1+b)
\]
with generic values. Scalar \(b\)-fusions give \(2 f_N^2=C^m\prod_{j<k\le m} D(x_k-x_j)^2\), nonzero, so the normalized columns specialize regularly. The right column \(w^i(v)\) fuses from the inner pairs out as a product of \(\Omega_{A_j B_j}\), subscripts here labeling partners of values \(x_j,x_j+b\). For \(w^i(-v)\) move the last site to the front and fuse from the outer pairs in; by cyclic covariance this gives a product of \(K_{B_j}(i)\Omega_{B_j A_j}\). Consequently \(Z=(1+y+y^{-1})^m\). In \(G\), a nonzero summand requires \(\alpha_{A_j}\ge\alpha_{B_j}\) at each partner pair by \(R_c(b+L)=0\); neutrality forces equality. Each pair then gives multiplier \(C\) for each of the three choices before the \(y\)-weight. Between any two pairs the remaining kernel factors multiply to 1 by \(R_s(d)^2 R_c(d+b)R_c(d-b)=1\). Thus \(G=U\) there also, proving the identity.

\paragraph{The homogeneous physical contraction.}
Finally set \(v=0,t=L\). Then \(T_\xi(L)^H\) as \(H\to\infty\) converges to a rank-one projection: in the path expansion with loops cancelled, through-lines to height \(H\) have vanishing weight by lifting to the planar honeycomb strip and using the strip asymptotic \(B_H\asymp H^{-1/4}\) from Section~\ref{c6:strip-asymptotics} (coordinate axes interchanged), and arch weights are summable by the same lifting and half-plane exit bound proved by rectangles in the strip argument. Indeed lift each strand separately starting from a fixed preimage of an endpoint, ignoring joint avoidance for an upper bound. Absolute weights factor by visited triangles (a two-arc tile at \(L\) visits two distinct ones), and seam and turn phases are unitary. Systems without through-lines hence converge to independent half-cylinder arch systems at the two ends. The empty entry is still 1.

This convergence makes 1 a simple eigenvalue of the physical row. Its normalized spectral projection is therefore analytic throughout a neighborhood of the homogeneous rapidities. Since \(f_N\) is the minimal common denominator of the normalized vector, \(f_N(0)\ne0\): any divisor component through that neighborhood would give genuine poles arbitrarily nearby. Thus the finite pairing can be evaluated at the physical point.

The sandwich defining \(Z\) is then the infinite-height limit of empty contractions with twist \(i\) below the horizontal line, twist \(-i\) above, and insertion of \(K(\delta)\)'s along segment \(AB\). A contractible oriented polygon, taken with interior on its left, gets phase \(i\) before the switch and insertion. Likewise an essential simple polygon with its lower side on its left, moving net towards the left, gets phase \(i\) before the switch. In either case, calling the left side \(D'\), the switch on the upward seam from \(A\) multiplies this by \((-1)^{1_{A\in D'}}\), and the insertion contributes \(\delta^{1_{A\in D'}-1_{B\in D'}}\) by oriented crossing count. Reversing orientation inverts the factor; the two sum to 0 unless the points are separated, when they sum to \(y+y^{-1}\). Tile weights without phases are precisely the honeycomb weights. For positive \(y\), exhaustion proves the stated \(\mathcal Z\) identity, and hence also as a Laurent polynomial (separating polygons bounded in number by disjoint crossings of the segment).
\end{proof}
